% Generated draft: requires mathematical and source review.
\documentclass[11pt]{article}
\usepackage[margin=25mm]{geometry}
\usepackage{fontspec}
\setmainfont{Latin Modern Roman}
\usepackage{amsmath,amssymb,mathtools,unicode-math}
\setmathfont{Latin Modern Math}
\usepackage{enumitem,needspace,xurl,hyperref}
\setlength{\parindent}{0pt}
\setlength{\parskip}{5pt}
\setlength{\emergencystretch}{3em}
\newcommand{\N}{\mathbb{N}}
\newcommand{\Z}{\mathbb{Z}}
\newcommand{\Q}{\mathbb{Q}}
\newcommand{\R}{\mathbb{R}}
\newcommand{\C}{\mathbb{C}}
\newcommand{\F}{\mathbb{F}}
\newcommand{\A}{\mathbb{A}}
\newcommand{\PP}{\mathbb{P}}
\newcommand{\E}{\mathbb{E}}
\newcommand{\eps}{\varepsilon}
\newcommand{\dd}{\,\mathrm{d}}
\begin{document}
\section*{UnsolvedMath ID: 30003442}
\textbf{Source problem code:} catalogue entry
\subsection{Definitions and mathematical statement}
Let \(m\) denote the number of variables and \(d\) the homogeneous degree. The supplied record does not specify their admissible ranges or any relation between them. A polynomial \(P(x_1,\ldots,x_m)\) is homogeneous of degree \(d\) if every monomial with nonzero coefficient has total degree \(d\).

Write \(D_{m,d}\) for the source-designated class of doubly stochastic stable homogeneous polynomials of degree \(d\) in \(m\) variables. The record does not specify the coefficient domain, the stability convention, or the normalization meant by ``doubly stochastic.'' These definitions must be recovered from the source before this class is mathematically determined.

For \(P\in D_{m,d}\), let \(\chi_P\) denote the mixed characteristic polynomial intended by the source, and let \(\lambda_1(P)\) denote its largest zero. No formula for \(\chi_P\) is supplied. The intended construction must also ensure that the largest zero exists and is real; those requirements cannot be checked from the record alone. Set
\[
P_\infty(x_1,\ldots,x_m)=\frac{(x_1+\cdots+x_m)^d}{m^d}.
\]
The supplied conjecture is that, for every admissible pair \((m,d)\),
\[
\forall P\in D_{m,d},\qquad
\lambda_1(P)\leq\lambda_1(P_\infty).
\]
This formulation retains the source's notation and normalization, but remains conditional on specifying the missing definitions, parameter restrictions, and conditions making both largest zeros meaningful.

The supplied dated literature assessment labels the problem partially solved: it reports that the bound is proved for symmetric stable doubly stochastic polynomials, while the unrestricted conjecture remains unresolved, and describes the symmetric case as a substantial but proper subclass. This is a supplied assessment, not an independently verified mathematical or literature conclusion.
\subsection{Short English statement}
The conjecture asks whether every polynomial in the intended doubly stochastic stable homogeneous class has largest mixed-characteristic-polynomial zero at most that of \((x_1+\cdots+x_m)^d/m^d\). The class normalization and mixed characteristic construction are missing, so the statement requires source clarification. Supplied dated curation reports a result for the symmetric subclass and labels the general problem partially solved.
\subsection{Sources}
\begin{itemize}
\item Supplied local problem registry record, permanent UnsolvedMath ID \(30003442\), title ``Largest Roots of Doubly Stochastic Stable Polynomials,'' collection ``Oberwolfach Reports Open Problems,'' exact source problem code \texttt{catalogue entry}. The record gives proposed year 2017 and names no individual proposer.
\item Supplied report citation: \emph{Real Algebraic Geometry With a View Toward Moment Problems and Optimization} (2018), Oberwolfach Reports, pp.~771--862, DOI \texttt{10.4171/owr/2017/14}. \url{https://doi.org/10.4171/owr/2017/14}. The supplied background describes Oberwolfach Reports as documenting mathematical workshops and their open-problem discussions.
\item Supplied literature reference: ``Branden, van der Waerden and Kadison--Singer with symmetries.'' \url{https://doi.org/10.4171/owr/2017/14}.
\item Supplied literature reference: ``Marcus--Spielman--Srivastava, mixed characteristic polynomials.'' \url{https://annals.math.princeton.edu/2015/182-1/p08}.
\item Literature checked at: \texttt{2026-08-21}. The partially solved label and the assessment concerning the symmetric subclass are supplied dated curation derived from the supplied OWR curation, not independently verified. The supplied background requests independent verification of time-sensitive status and cited claims.
\end{itemize}
\end{document}
